\chapter{提案システムの設計と実装}
\label{chap:system}

本章では，複数のトレイルカメラで取得した画像をエッジサーバに集約し，動物が含まれる可能性のある画像をクラウドへ転送するシステムについて述べる．
説明の対象はICOIACT用原稿に記載した試作構成とし，カメラ，エッジサーバおよびクラウドサーバの機能と処理手順を整理する．
一次選別にはCOCOで事前学習されたYOLOv8n，再判定にはMegaDetector v5aを用いる．

\section{システムの要件}
農地や山間部で複数地点を監視する場合，撮影地点ごとの電源確保と，遠隔地へ情報を送る通信経路の確保が必要となる．
さらに，PIRセンサの反応で発生する空撮影画像をすべて転送すると，広域通信のデータ量と利用者の確認作業が増加する．
第\ref{chap:related}章の関連研究を踏まえ，本システムでは次の事項を設計要件とする．
\begin{enumerate}
  \item 複数地点のカメラを一つのエッジサーバへ接続し，広域通信回線を集約できること．
  \item カメラは検知，撮影および局所的な画像転送を担い，撮影のない時間は休止できること．
  \item 広域通信の前に画像を選別し，動物の見逃しを考慮しながら不要な転送を減らせること．
  \item 転送画像を遠隔側で再判定し，動物の検出結果と画像を利用者へ通知できること．
  \item クラウドが受信した画像と解析結果を保存し，事後の確認に利用できること．
\end{enumerate}
これらは設計上の要件であり，長期間の自律稼働や通知時間の達成を保証するものではない．
通信範囲，接続台数，撮影頻度および電源条件によって運用可能な範囲は変化するため，第\ref{chap:evaluation}章で評価条件と対応付けて検討する．
\writehere{目標監視範囲，想定カメラ台数，許容通知時間，連続稼働期間および費用上限の数値を，修士論文の評価対象に合わせて確定する．}

\section{システム全体構成}
提案システムは，エッジカメラ，エッジサーバおよびクラウドサーバの3層で構成する．
図\ref{fig:system-overview}に，各層の役割と通信経路を示す．
カメラは動体検知を契機に画像を取得し，ローカルネットワークを介してエッジサーバへ送信する．
エッジサーバは複数カメラからの画像を受信し，YOLOv8nで一次選別して，選別を通過した画像をLTE回線からクラウドへ送信する．
クラウドサーバはMegaDetector v5aで再判定し，受信画像を保存するとともに，動物を検出した場合に画像付きの電子メールで通知する．

\begin{figure}[tbp]
  \centering
  \fbox{\begin{minipage}{0.88\linewidth}
    \centering
    \textbf{複数のエッジカメラ}\\
    PIR検知・3枚連続撮影・画像送信・待機\\[1mm]
    $\downarrow$\quad 2.4\,GHz Wi-Fi／HTTP\\[1mm]
    \textbf{共有エッジサーバ（Raspberry Pi 4B）}\\
    画像集約・YOLOv8nによる一次選別\\[1mm]
    $\downarrow$\quad LTE／HTTP（選別を通過した画像）\\[1mm]
    \textbf{クラウドサーバ}\\
    MegaDetector v5aによる再判定・画像と結果の保存\\[1mm]
    $\downarrow$\quad 動物検出時に画像付きメール\\[1mm]
    \textbf{利用者}
  \end{minipage}}
  \caption{提案システムの構成と画像・通知の流れ}
  \label{fig:system-overview}
\end{figure}

表\ref{tab:system-components}に試作構成の主要な機器とモデルを示す．
この構成では，各カメラに推論用の計算機やLTE回線を配置せず，それらをエッジサーバ側で共有する．
Satoら\cite{Sato2022}の親機・子機への機能分離を踏まえ，カメラ側を画像取得に特化させるとともに，共有側での一次選別と遠隔側での再判定を組み合わせる．
ただし，カメラからエッジサーバへの通信は一次選別前に行うため，本方式が直接削減するのは主にエッジサーバからクラウドへの画像転送である．

\begin{table}[tbp]
  \centering
  \small
  \caption{ICOIACT用原稿に基づく試作構成}
  \label{tab:system-components}
  \begin{tabularx}{\linewidth}{p{0.23\linewidth}XX}
    \toprule
    構成要素 & 機器・モデル & 役割 \\
    \midrule
    カメラ基板 & XIAO ESP32S3 Sense & 撮影とWi-Fi通信 \\
    動体センサ & Panasonic PaPIRs & 撮影開始のトリガー \\
    エッジサーバ & Raspberry Pi 4B（メモリ4\,GB） & 画像受信，一次選別，転送 \\
    通信機器 & NEC Aterm HT100LN & Wi-FiアクセスポイントとLTE接続 \\
    一次選別モデル & YOLOv8n（COCO事前学習） & 動物候補を含む画像群の選別 \\
    再判定モデル & MegaDetector v5a & クラウドでの動物検出 \\
    \bottomrule
  \end{tabularx}
\end{table}

\section{エッジカメラ}
\subsection{ハードウェア構成}
エッジカメラには，Seeed StudioのXIAO ESP32S3 Senseを用いる．
カメラ機能とWi-Fi通信機能を備える基板を用い，PIRセンサの信号を契機に撮影して，画像をエッジサーバへ転送する．
動体検知にはPanasonic PaPIRsを用いる．
PIRセンサの反応は動物の存在を確定するものではないため，カメラ側では動物種の識別や画像選別を行わず，画像の内容はエッジサーバ以降で判定する．

撮影地点での電池駆動を想定し，待機にはESP32-S3のDeep-sleepを利用する．
ただし，マイコンの休止だけで装置全体の消費電力が決まるわけではなく，周辺回路を含めた待機時と動作時の電力を評価する必要がある．
ICOIACT用原稿には単3形電池8本，電源電圧6\,Vの構成が記載されているが，電池の接続方法と電源回路は原稿だけでは確定できない．
また，消費電流の測定条件は供給電圧5.0\,Vであり，電池構成の記述と測定時の条件は区別して扱う．
\writehere{PIRセンサの型番，カメラモジュール，配線，記憶媒体，電池の直並列接続と電源回路，筐体の防水構造，夜間照明の有無と制御を実機資料から補う．}

\subsection{撮影・保存・送信処理}
カメラの処理は，待機，復帰，連続撮影，画像送信および待機への復帰からなる．
待機中はDeep-sleepを用い，PIRセンサの検知により起動して3枚の画像を連続撮影する．
その後，Wi-Fi経由のHTTP通信で画像をエッジサーバへ送り，送信後に所定の間隔を置いて待機状態へ戻る．
本章では，一度の検知に伴う3枚の撮影とその送信処理を撮影サイクルと呼ぶ．

連続撮影によって，同じ検知の前後で姿勢や画面内の位置が異なる画像を取得できる．
後述する一次選別では，この画像群に対する判定をまとめて転送の可否を決める．
ただし，3枚がそれぞれ独立した動物の出現を表すとは限らず，画像枚数と出現回数は区別する．
\writehere{画像解像度，圧縮形式・品質，撮影間隔，待機へ戻るまでの時間，送信前の保存方法と保存期間を，評価に用いたファームウェアで確認する．}

\section{エッジサーバ}
\subsection{画像の受信と管理}
エッジサーバには，メモリ4\,GBのRaspberry Pi 4Bを用いる．
同サーバはカメラからのHTTP送信を受け付け，各画像を一次選別の処理へ渡す．
選別を通過した画像は保存用ディレクトリへ移し，クラウドへの転送対象とする．
カメラ側を間欠動作とする一方，エッジサーバは画像の到着を待ち受ける構成であり，推論を実行していない時間にも受信機能を維持する．

複数カメラの画像を撮影サイクルごとに扱うには，カメラの識別情報，サイクルの識別情報およびサイクル内の画像順序を対応付ける必要がある．
また，短時間に画像が集中した場合には受信と推論の待ち時間が生じるため，処理の順序や保存領域の容量は通知遅延と運用に関わる．
\writehere{OS・ライブラリの版，画像識別子とファイル名，サイクルの対応付け，受信待ち行列，同時接続数，保存期間を評価対象の実装に合わせて記載する．}

\subsection{一次選別とクラウド転送}
一次選別には，COCOデータセットで事前学習されたYOLOv8nを用いる．
計算資源の限られるエッジサーバで画像を選別するため，小規模なモデルを採用する．
その目的は動物種を厳密に同定することではなく，動物が写っている可能性のある撮影サイクルを残し，明らかな空撮影のクラウド転送を減らすことである．
ICOIACT用原稿の評価では，動物画像を誤って除外することを抑えるため，信頼度閾値を$\tau_{\mathrm{edge}}=0.05$としている．

撮影サイクル$k$の3枚の画像を$x_{k,1},x_{k,2},x_{k,3}$とする．
画像$x$について，一次選別の対象とする動物クラスの検出候補があり，その最大信頼度が$\tau_{\mathrm{edge}}$以上の場合に$d_{\mathrm{edge}}(x)=1$，それ以外の場合に0とする．
サイクルの転送判定$D_k$は，次式で表せる．
\begin{equation}
  D_k=\max_{j\in\{1,2,3\}}d_{\mathrm{edge}}(x_{k,j})．
  \label{eq:system-cycle-selection}
\end{equation}
すなわち，3枚のうち少なくとも1枚で動物候補を検出すれば，当該サイクルの画像群を転送対象とする．
いずれの画像でも検出しなかった場合はクラウドへ転送しない．
これは画像群を単位とする選別であり，転送対象のサイクル内に空撮影画像が含まれる可能性は残る．

一次選別の閾値を低くすると，誤検出を含むサイクルが転送される可能性も高まる．
本システムでは，そのような候補をクラウドで再判定する一方，一次選別で除外した画像をクラウド側の処理によって回復することはできない．
したがって，Cunhaら\cite{Cunha2021}の評価の観点も踏まえ，削減率だけでなく動物画像の保持と推論時間を併せて確認する．
\writehere{対象クラスの一覧，モデル重みと入力サイズ，サイクル通過時の実際の転送枚数，画像欠落時の扱いを確認する．式はICOIACT原稿の処理説明を表すものであり，第5章の集計コードについては画像単位とサイクル単位の対応を別途照合する．}

\section{クラウドサーバ}
クラウドサーバは，エッジサーバから転送された画像に対してMegaDetector v5aによる再判定を行う．
一次選別と異なる検出モデルを用いて動物の有無を調べ，動物を検出した場合には画像を添付した電子メールで利用者へ通知する．
MegaDetectorの動物カテゴリは種名を表すものではないため，ここでの通知条件は動物の検出とし，特定の獣種を識別した通知とは区別する．

受信した画像と解析結果は，動物の検出の有無にかかわらずクラウド側に蓄積する．
動物を検出しなかった場合には，通知せず保存のみを行う．
この保存対象は一次選別を通過してクラウドが受信した画像であり，カメラの全撮影画像がクラウドに残るわけではない．
保存画像は，判定結果の事後確認やモデル改善に向けた分析に利用できる．

再判定は不要な通知を抑えるための処理であるが，MegaDetectorにも誤検出と見逃しがある\cite{Fennell2022}．
そのため，クラウドでの検出結果を正解ラベルと同一視せず，人手確認による評価と，モデル出力との一致を測る評価とを分ける．
\writehere{クラウドの実行環境，再判定の閾値，画像単位とサイクル単位の通知条件，連続検出時の通知のまとめ方，保存先と保存期間を確認する．閲覧画面については実装を確認してから記載する．}

\section{通信と障害時の処理}
\subsection{通信経路}
カメラからエッジサーバまでのローカル通信には，NEC Aterm HT100LNの2.4\,GHz帯Wi-Fiアクセスポイントを利用する．
各カメラはこのネットワークへ接続し，HTTPで画像を送る．
エッジサーバからクラウドサーバまでは，同機器のLTE回線を介してインターネットへ接続し，HTTPで選別後の画像を転送する．
この構成では，広域通信の契約をカメラごとに用意する必要がない一方，カメラの設置範囲はローカル通信の到達性によって制約される．

第\ref{sec:background-halow}節で説明したWi-Fi HaLowは通信範囲の拡張を検討する際の技術候補であり，本章の試作構成で使用した無線方式とは区別する．
HaLowを用いる場合には対応機器を含めた接続構成を定義し，通信成功率，転送時間および電力量を別途評価する必要がある．

\subsection{障害時の検討事項}
処理経路には，カメラからエッジサーバへの転送，エッジでの推論，クラウドへの転送，再判定およびメール通知が含まれる．
いずれかの段階で失敗すると，撮影できても利用者へ通知されない場合がある．
特に，広域回線を共有する構成では，エッジサーバまたはLTE回線の停止が，接続する複数カメラの通知に影響する．

通信失敗時に画像を保持して再送する方法，受信済み画像の重複処理を防ぐ方法，保存領域の容量不足に対応する方法は，継続運用のために整理すべき事項である．
ICOIACT用原稿ではこれらの具体的な処理条件を示していないため，本章では実装済み機能として断定しない．
\writehere{評価対象のコード版を固定し，接続・応答のタイムアウト，再送回数，送信成功の判定，通信断時の保存，重複受信の扱い，容量不足時の動作を確認して記載する．}
